% 1.6
% Ficheiro incluído por main.tex

\section[Firewall do Windows: permitir ping do Ubuntu]{Firewall do Windows: permitir ping do Ubuntu}

\passo{Passo 1 - Nova regra de entrada}

\emph{Iniciar} → escrever \textbf{Windows Defender Firewall with Advanced Security} e abrir. Painel esquerdo: \textbf{Inbound Rules}; painel direito: \textbf{New Rule...}

\passo{Passo 2 - Tipo de regra e protocolo}

\textbf{Custom} → Next → \textbf{All programs} → Next → \textbf{Protocol type: ICMPv4} → Next.

\passo{Passo 3 - IP de origem}

Em \emph{Which remote IP addresses does this rule apply to?} → \textbf{These IP addresses} → \textbf{Add...} → \texttt{\textless{}IP\_CLIENTE\textgreater{}} → OK. Repetir \textbf{Add...} com \texttt{\textless{}IP\_CLIENTE\_WB\textgreater{}} (o cliente Winbind de 2.2) → OK → Next.

\passo{Passo 4 - Permitir e dar nome}

\textbf{Allow the connection} → Next → Next → Nome: \textbf{\texttt{SR2 Allow Ping Ubuntu}} → \textbf{Finish}.

\emph{Verificação esperada:} a regra aparece ativa (visto verde) em \emph{Inbound Rules}. O teste de ping faz-se a partir do Ubuntu (2.1.2, Passo 11).

